\clearpage\thispagestyle{empty}
\pdfbookmark[0]{Quellen zu Komet (Briefe vom 8. und 16. September 1811)}{comet1811septletters-sources}
\begin{center}\textbf{Quellen zu Komet (Briefe vom 8. und 16. September 1811)}\end{center}
\noindent Diese KI-Übertragung ins heutige Deutsch folgt dem Bericht und den Auszügen aus den Briefen von Carl Friedrich Gauß in \href{https://gdz.sub.uni-goettingen.de/id/PPN236007467}{Gauß’ \textit{Werken}, Bd.~VI} (Göttingen, 1874), S.~335--337. Die Auszüge erschienen in Beitrag XXXI, \textit{Ueber den grossen Cometen von 1811}, in der \textit{Monatlichen Correspondenz}, Bd.~XXIV, September 1811, S.~\href{https://api.digitale-sammlungen.de/iiif/image/v2/bsb10538613_00314/full/1600,/0/default.jpg}{304}, \href{https://api.digitale-sammlungen.de/iiif/image/v2/bsb10538613_00315/full/1600,/0/default.jpg}{305}, \href{https://api.digitale-sammlungen.de/iiif/image/v2/bsb10538613_00316/full/1600,/0/default.jpg}{306} und \href{https://api.digitale-sammlungen.de/iiif/image/v2/bsb10538613_00317/full/1600,/0/default.jpg}{307}. Das digitalisierte Exemplar stammt aus der Bayerischen Staatsbibliothek.
\par\smallskip\noindent Die Briefdaten stehen in der Zeitschrift in Fußnoten auf S.~304 und 306, in der Werkausgabe im Text. Der Erstdruck und die Werkausgabe geben für den 17. September die Stunde 0 an. Der Punkt vor der 27 in der Septemberephemeride steht in der Werkausgabe. Die historische Fußnote verweist auf das Augustheft, S.~180. KI-Übertragung und KI-Satz.
\clearpage
